%%%%%%%%%%%%%%%%%%%%%%%%%%%%%%%%%%%%%%%%%%%%%%%%%%%%%%%%%%%%
% 0.9 CARDINALITY AND INFINITY
% The Composition of Advanced Mathematics
%%%%%%%%%%%%%%%%%%%%%%%%%%%%%%%%%%%%%%%%%%%%%%%%%%%%%%%%%%%%

\section{Cardinality and Infinity}
\label{sec:cardinality-infinity}

\prerequisite{
Familiarity with sets from \cref{sec:elementary-set-theory}, functions from
\cref{sec:functions}, and proof techniques from
\cref{sec:methods-of-proof}.
}

\learningobjectives{
By the end of this section, the reader should be able to:
\begin{itemize}
    \item interpret cardinality as the size of a set;
    \item compare finite sets using bijections;
    \item distinguish finite, infinite, countably infinite, countable, and uncountable sets;
    \item understand how infinite sets can be placed in bijection with proper subsets of themselves;
    \item construct bijections involving \(\N\), \(\Z\), and subsets of these sets;
    \item understand why \(\N\times\N\) is countable;
    \item understand why countable unions of countable sets are countable;
    \item explain why \(\Q\) is countable;
    \item distinguish density from cardinality;
    \item understand Cantor's diagonal argument;
    \item explain why \(\R\) is uncountable;
    \item compare infinite cardinalities;
    \item interpret the cardinal \(\aleph_0\);
    \item understand the continuum cardinality \(\mathfrak c\);
    \item state the Cantor--Bernstein--Schröder theorem;
    \item understand Cantor's theorem for power sets;
    \item recognize that there is no largest infinity.
\end{itemize}
}

%%%%%%%%%%%%%%%%%%%%%%%%%%%%%%%%%%%%%%%%%%%%%%%%%%%%%%%%%%%%
\subsection{What Does It Mean for a Set to Have a Size?}
\label{subsec:meaning-cardinality}
%%%%%%%%%%%%%%%%%%%%%%%%%%%%%%%%%%%%%%%%%%%%%%%%%%%%%%%%%%%%

For a finite set, the idea of size is familiar.

If

\[ A=\{a,b,c,d\}, \]

then \(A\) contains four elements.

We write

\[ |A|=4. \]

The notation

\[ |A| \]

is read

\begin{center}
``the cardinality of \(A\)''
\end{center}

or

\begin{center}
``the number of elements of \(A\).''
\end{center}

For finite sets, cardinality simply counts the elements.

The concept becomes much more interesting for infinite sets.

%%%%%%%%%%%%%%%%%%%%%%%%%%%%%%%%%%%%%%%%%%%%%%%%%%%%%%%%%%%%
\subsection{Cardinality}
\label{subsec:cardinality-definition}
%%%%%%%%%%%%%%%%%%%%%%%%%%%%%%%%%%%%%%%%%%%%%%%%%%%%%%%%%%%%

\begin{definitionbox}{Cardinality}
The \emph{cardinality} of a set is a mathematical description of the size of
that set.

For finite sets, cardinality is the number of elements in the set.
\end{definitionbox}

For example,

\[ A=\{2,4,6,8,10\} \]

has cardinality

\[ |A|=5. \]

Similarly,

\[ |\varnothing|=0. \]

%%%%%%%%%%%%%%%%%%%%%%%%%%%%%%%%%%%%%%%%%%%%%%%%%%%%%%%%%%%%
\subsection{Comparing Sets Without Counting}
\label{subsec:comparing-without-counting}
%%%%%%%%%%%%%%%%%%%%%%%%%%%%%%%%%%%%%%%%%%%%%%%%%%%%%%%%%%%%

Suppose one has two collections of objects and wants to determine whether they
have the same size.

Instead of counting both collections, one may pair each object from the first
collection with exactly one object from the second.

If every object is paired once and no object is left over, the two collections
have the same size.

This idea extends naturally to infinite sets.

%%%%%%%%%%%%%%%%%%%%%%%%%%%%%%%%%%%%%%%%%%%%%%%%%%%%%%%%%%%%
\subsection{Equal Cardinality}
\label{subsec:equal-cardinality}
%%%%%%%%%%%%%%%%%%%%%%%%%%%%%%%%%%%%%%%%%%%%%%%%%%%%%%%%%%%%

\begin{definitionbox}{Equal Cardinality}
Two sets \(A\) and \(B\) have the same cardinality if there exists a bijection
\[ f:A\to B. \]
\end{definitionbox}

We may write

\[ |A|=|B|. \]

Thus cardinality is fundamentally connected to bijective functions.

%%%%%%%%%%%%%%%%%%%%%%%%%%%%%%%%%%%%%%%%%%%%%%%%%%%%%%%%%%%%
\subsection{Worked Example: Equal Finite Cardinality}
\label{subsec:worked-equal-finite-cardinality}
%%%%%%%%%%%%%%%%%%%%%%%%%%%%%%%%%%%%%%%%%%%%%%%%%%%%%%%%%%%%

\begin{workedexamplebox}{Compare Two Finite Sets}

Let

\[ A=\{a,b,c\},\qquad B=\{4,7,9\}. \]

Define

\[ f(a)=4,\qquad f(b)=7,\qquad f(c)=9. \]

\begin{tabularx}{\textwidth}{@{}>{\raggedright\arraybackslash}p{0.42\textwidth}X@{}}
Each element of \(A\) receives one output. & \(f\) is a function\\
Distinct elements receive distinct outputs. & \(f\) is injective\\
Every element of \(B\) is used. & \(f\) is surjective
\end{tabularx}

Therefore \(f\) is bijective.

\begin{answerbox}
\[ |A|=|B|=3 \]
\end{answerbox}
\end{workedexamplebox}

%%%%%%%%%%%%%%%%%%%%%%%%%%%%%%%%%%%%%%%%%%%%%%%%%%%%%%%%%%%%
\subsection{Finite Sets}
\label{subsec:finite-sets}
%%%%%%%%%%%%%%%%%%%%%%%%%%%%%%%%%%%%%%%%%%%%%%%%%%%%%%%%%%%%

\begin{definitionbox}{Finite Set}
A set \(A\) is \emph{finite} if either \(A=\varnothing\) or there exists some
positive integer \(n\) such that
\[ |A|=n. \]
\end{definitionbox}

Equivalently, a nonempty finite set can be placed in bijection with

\[ \{1,2,\ldots,n\} \]

for some positive integer \(n\).

%%%%%%%%%%%%%%%%%%%%%%%%%%%%%%%%%%%%%%%%%%%%%%%%%%%%%%%%%%%%
\subsection{Infinite Sets}
\label{subsec:infinite-sets}
%%%%%%%%%%%%%%%%%%%%%%%%%%%%%%%%%%%%%%%%%%%%%%%%%%%%%%%%%%%%

\begin{definitionbox}{Infinite Set}
A set is \emph{infinite} if it is not finite.
\end{definitionbox}

Examples include

\[ \N,\qquad \Z,\qquad \Q,\qquad \R,\qquad \C. \]

However, these infinite sets do not all have the same cardinality.

Understanding this fact is one of the central ideas of set theory.

%%%%%%%%%%%%%%%%%%%%%%%%%%%%%%%%%%%%%%%%%%%%%%%%%%%%%%%%%%%%
\subsection{A Strange Property of Infinite Sets}
\label{subsec:infinite-proper-subset}
%%%%%%%%%%%%%%%%%%%%%%%%%%%%%%%%%%%%%%%%%%%%%%%%%%%%%%%%%%%%

For finite sets, a proper subset always has strictly smaller cardinality.

For example,

\[ \{1,2\}\subsetneq\{1,2,3\} \]

and

\[ 2<3. \]

Infinite sets behave differently.

Consider

\[ \N=\{1,2,3,\ldots\} \]

and the set of positive even integers

\[ E=\{2,4,6,\ldots\}. \]

Clearly,

\[ E\subsetneq\N. \]

Yet the function

\[ f:\N\to E,\qquad f(n)=2n \]

is a bijection.

Therefore,

\[ |\N|=|E|. \]

An infinite set can have the same cardinality as one of its proper subsets.

%%%%%%%%%%%%%%%%%%%%%%%%%%%%%%%%%%%%%%%%%%%%%%%%%%%%%%%%%%%%
\subsection{Dedekind-Infinite Behavior}
\label{subsec:dedekind-infinite-preview}
%%%%%%%%%%%%%%%%%%%%%%%%%%%%%%%%%%%%%%%%%%%%%%%%%%%%%%%%%%%%

This phenomenon can be used as one characterization of infinite sets in
standard set theory.

A set may be infinite while still being pairable one-to-one with a proper
subset of itself.

This is fundamentally different from finite sets.

%%%%%%%%%%%%%%%%%%%%%%%%%%%%%%%%%%%%%%%%%%%%%%%%%%%%%%%%%%%%
\subsection{Countably Infinite Sets}
\label{subsec:countably-infinite}
%%%%%%%%%%%%%%%%%%%%%%%%%%%%%%%%%%%%%%%%%%%%%%%%%%%%%%%%%%%%

\begin{definitionbox}{Countably Infinite Set}
A set \(A\) is \emph{countably infinite} if there exists a bijection
\[ f:\N\to A. \]
\end{definitionbox}

Intuitively, the elements of \(A\) can be arranged in an infinite list:

\[ a_1,a_2,a_3,\ldots. \]

Every element must appear somewhere on the list.

%%%%%%%%%%%%%%%%%%%%%%%%%%%%%%%%%%%%%%%%%%%%%%%%%%%%%%%%%%%%
\subsection{The Cardinal \(\aleph_0\)}
\label{subsec:aleph-zero}
%%%%%%%%%%%%%%%%%%%%%%%%%%%%%%%%%%%%%%%%%%%%%%%%%%%%%%%%%%%%

The symbol

\[ \aleph_0 \]

uses the Hebrew letter \emph{aleph}.

It is pronounced

\begin{center}
``aleph-null''
\end{center}

or

\begin{center}
``aleph-zero.''
\end{center}

The subscript \(0\) is read ``zero'' or sometimes ``naught.''

\begin{definitionbox}{Aleph-Zero}
The cardinality of the natural numbers is denoted
\[ |\N|=\aleph_0. \]
\end{definitionbox}

The cardinal \(\aleph_0\) represents the size of every countably infinite set.

%%%%%%%%%%%%%%%%%%%%%%%%%%%%%%%%%%%%%%%%%%%%%%%%%%%%%%%%%%%%
\subsection{Countable Sets}
\label{subsec:countable-sets}
%%%%%%%%%%%%%%%%%%%%%%%%%%%%%%%%%%%%%%%%%%%%%%%%%%%%%%%%%%%%

The word \emph{countable} is commonly used slightly more broadly.

\begin{definitionbox}{Countable Set}
A set is \emph{countable} if it is finite or countably infinite.
\end{definitionbox}

Thus every finite set is countable, and every set with cardinality
\(\aleph_0\) is countable.

%%%%%%%%%%%%%%%%%%%%%%%%%%%%%%%%%%%%%%%%%%%%%%%%%%%%%%%%%%%%
\subsection{Infinite Subsets of Countable Sets}
\label{subsec:infinite-subsets-countable}
%%%%%%%%%%%%%%%%%%%%%%%%%%%%%%%%%%%%%%%%%%%%%%%%%%%%%%%%%%%%

An infinite subset of a countable set is itself countably infinite.

For example, the prime numbers

\[ \{2,3,5,7,11,\ldots\} \]

form an infinite subset of \(\N\).

Therefore the prime numbers are countably infinite.

This does not mean there are ``as many primes as integers'' in terms of
density or frequency.

It means there is a bijection between the two sets.

%%%%%%%%%%%%%%%%%%%%%%%%%%%%%%%%%%%%%%%%%%%%%%%%%%%%%%%%%%%%
\subsection{The Integers Are Countable}
\label{subsec:integers-countable}
%%%%%%%%%%%%%%%%%%%%%%%%%%%%%%%%%%%%%%%%%%%%%%%%%%%%%%%%%%%%

At first, \(\Z\) may seem larger than \(\N\) because it contains positive
integers, negative integers, and zero.

Nevertheless,

\[ |\Z|=|\N|. \]

One way to list the integers is

\[ 0,1,-1,2,-2,3,-3,\ldots. \]

Every integer eventually appears in this list.

Therefore \(\Z\) is countably infinite.

%%%%%%%%%%%%%%%%%%%%%%%%%%%%%%%%%%%%%%%%%%%%%%%%%%%%%%%%%%%%
\subsection{A Bijection from \(\N\) to \(\Z\)}
\label{subsec:bijection-n-z}
%%%%%%%%%%%%%%%%%%%%%%%%%%%%%%%%%%%%%%%%%%%%%%%%%%%%%%%%%%%%

A formal bijection may be defined by

\[ f(n)=
\begin{cases}
\dfrac n2,&n\text{ even},\\[4pt]
-\dfrac{n-1}{2},&n\text{ odd}.
\end{cases} \]

With \(\N=\{1,2,3,\ldots\}\), this gives

\[ f(1)=0,\quad f(2)=1,\quad f(3)=-1,\quad f(4)=2,\quad f(5)=-2,\ldots \]

Every integer occurs exactly once.

%%%%%%%%%%%%%%%%%%%%%%%%%%%%%%%%%%%%%%%%%%%%%%%%%%%%%%%%%%%%
\subsection{Worked Example: Positive Multiples of Three}
\label{subsec:worked-multiples-three}
%%%%%%%%%%%%%%%%%%%%%%%%%%%%%%%%%%%%%%%%%%%%%%%%%%%%%%%%%%%%

\begin{workedexamplebox}{Show a Proper Subset of \(\N\) Has the Same Cardinality}

Let

\[ M=\{3,6,9,12,\ldots\}. \]

Define

\[ f:\N\to M,\qquad f(n)=3n. \]

\begin{tabularx}{\textwidth}{@{}>{\raggedright\arraybackslash}p{0.42\textwidth}X@{}}
Suppose \(f(m)=f(n)\). & \(3m=3n\Rightarrow m=n\)\\
Thus \(f\) is injective. & Distinct inputs have distinct outputs\\
Every element of \(M\) has form \(3n\). & \(f\) is surjective
\end{tabularx}

\begin{answerbox}
\[ |\N|=|M|=\aleph_0 \]
\end{answerbox}
\end{workedexamplebox}

%%%%%%%%%%%%%%%%%%%%%%%%%%%%%%%%%%%%%%%%%%%%%%%%%%%%%%%%%%%%
\subsection{Cartesian Products and Countability}
\label{subsec:cartesian-countability}
%%%%%%%%%%%%%%%%%%%%%%%%%%%%%%%%%%%%%%%%%%%%%%%%%%%%%%%%%%%%

Consider

\[ \N\times\N. \]

Its elements are ordered pairs

\[ (m,n) \]

of natural numbers.

It may initially appear that two coordinates should produce a larger
infinity.

Surprisingly,

\[ |\N\times\N|=\aleph_0. \]

Thus \(\N\times\N\) is countably infinite.

%%%%%%%%%%%%%%%%%%%%%%%%%%%%%%%%%%%%%%%%%%%%%%%%%%%%%%%%%%%%
\subsection{Diagonal Enumeration of \(\N\times\N\)}
\label{subsec:diagonal-nxn}
%%%%%%%%%%%%%%%%%%%%%%%%%%%%%%%%%%%%%%%%%%%%%%%%%%%%%%%%%%%%

The ordered pairs can be listed diagonally.

For example:

\[ (1,1), \]

then

\[ (1,2),(2,1), \]

then

\[ (1,3),(2,2),(3,1), \]

then

\[ (1,4),(2,3),(3,2),(4,1), \]

and so on.

Every ordered pair appears on exactly one diagonal corresponding to a fixed
value of

\[ m+n. \]

Thus all pairs can be placed into a single sequence.

%%%%%%%%%%%%%%%%%%%%%%%%%%%%%%%%%%%%%%%%%%%%%%%%%%%%%%%%%%%%
\subsection{Finite Cartesian Powers of \(\N\)}
\label{subsec:finite-cartesian-powers-n}
%%%%%%%%%%%%%%%%%%%%%%%%%%%%%%%%%%%%%%%%%%%%%%%%%%%%%%%%%%%%

The same idea extends to higher finite dimensions.

For every positive integer \(k\),

\[ |\N^k|=\aleph_0. \]

Thus

\[ \N,\qquad \N^2,\qquad \N^3,\qquad \ldots \]

all have the same cardinality.

This is another example of how infinite cardinality differs dramatically from
finite size.

%%%%%%%%%%%%%%%%%%%%%%%%%%%%%%%%%%%%%%%%%%%%%%%%%%%%%%%%%%%%
\subsection{Countable Unions}
\label{subsec:countable-unions}
%%%%%%%%%%%%%%%%%%%%%%%%%%%%%%%%%%%%%%%%%%%%%%%%%%%%%%%%%%%%

Suppose

\[ A_1,A_2,A_3,\ldots \]

are countable sets.

Then their union

\[ \bigcup_{n=1}^{\infty}A_n \]

is countable.

This principle is often summarized as:

\begin{importantbox}
A countable union of countable sets is countable.
\end{importantbox}

This result is extremely useful when proving that complicated sets are still
countable.

%%%%%%%%%%%%%%%%%%%%%%%%%%%%%%%%%%%%%%%%%%%%%%%%%%%%%%%%%%%%
\subsection{Why Countable Unions Remain Countable}
\label{subsec:why-countable-unions}
%%%%%%%%%%%%%%%%%%%%%%%%%%%%%%%%%%%%%%%%%%%%%%%%%%%%%%%%%%%%

Imagine arranging each countable set as a row:

\[
\begin{array}{cccc}
a_{11}&a_{12}&a_{13}&\cdots\\
a_{21}&a_{22}&a_{23}&\cdots\\
a_{31}&a_{32}&a_{33}&\cdots\\
\vdots&\vdots&\vdots&
\end{array}
\]

The entries may be traversed diagonally in the same way as
\(\N\times\N\).

If duplicates occur, simply omit repeated entries.

This produces a single list containing every element of the union.

%%%%%%%%%%%%%%%%%%%%%%%%%%%%%%%%%%%%%%%%%%%%%%%%%%%%%%%%%%%%
\subsection{The Rational Numbers}
\label{subsec:rational-countability}
%%%%%%%%%%%%%%%%%%%%%%%%%%%%%%%%%%%%%%%%%%%%%%%%%%%%%%%%%%%%

Recall that every rational number can be written as

\[ \frac pq \]

where

\[ p\in\Z,\qquad q\in\N. \]

Thus the rational numbers arise from pairs

\[ (p,q)\in\Z\times\N. \]

Since \(\Z\) and \(\N\) are countable, their Cartesian product is countable.

This suggests that \(\Q\) is countable.

%%%%%%%%%%%%%%%%%%%%%%%%%%%%%%%%%%%%%%%%%%%%%%%%%%%%%%%%%%%%
\subsection{The Rational Numbers Are Countable}
\label{subsec:q-countable}
%%%%%%%%%%%%%%%%%%%%%%%%%%%%%%%%%%%%%%%%%%%%%%%%%%%%%%%%%%%%

Every rational number is represented by at least one pair

\[ (p,q)\in\Z\times\N. \]

Since

\[ |\Z\times\N|=\aleph_0, \]

the collection of such pairs is countable.

Different pairs may represent the same rational number:

\[ \frac12=\frac24=\frac36. \]

Removing these duplicates does not make the set uncountable.

Therefore,

\[ |\Q|=\aleph_0. \]

%%%%%%%%%%%%%%%%%%%%%%%%%%%%%%%%%%%%%%%%%%%%%%%%%%%%%%%%%%%%
\subsection{A Surprising Consequence}
\label{subsec:q-surprising}
%%%%%%%%%%%%%%%%%%%%%%%%%%%%%%%%%%%%%%%%%%%%%%%%%%%%%%%%%%%%

The rational numbers are distributed throughout the real number line.

Between any two distinct real numbers there exists a rational number.

Nevertheless,

\[ |\Q|=|\N|. \]

Thus geometric density and cardinality are different concepts.

%%%%%%%%%%%%%%%%%%%%%%%%%%%%%%%%%%%%%%%%%%%%%%%%%%%%%%%%%%%%
\subsection{Density versus Cardinality}
\label{subsec:density-vs-cardinality}
%%%%%%%%%%%%%%%%%%%%%%%%%%%%%%%%%%%%%%%%%%%%%%%%%%%%%%%%%%%%

\begin{definitionbox}{Dense Subset}
A subset \(D\subseteq\R\) is \emph{dense in \(\R\)} if every open interval
contains an element of \(D\).
\end{definitionbox}

The rational numbers are dense in \(\R\).

The irrational numbers are also dense in \(\R\).

Density describes how a set is distributed.

Cardinality describes how many elements it contains.

These ideas should not be confused.

%%%%%%%%%%%%%%%%%%%%%%%%%%%%%%%%%%%%%%%%%%%%%%%%%%%%%%%%%%%%
\subsection{Uncountable Sets}
\label{subsec:uncountable-sets}
%%%%%%%%%%%%%%%%%%%%%%%%%%%%%%%%%%%%%%%%%%%%%%%%%%%%%%%%%%%%

\begin{definitionbox}{Uncountable Set}
A set is \emph{uncountable} if it is not countable.
\end{definitionbox}

An uncountable set cannot be placed into a finite or infinite sequence that
lists every element.

The real numbers provide the most important introductory example.

%%%%%%%%%%%%%%%%%%%%%%%%%%%%%%%%%%%%%%%%%%%%%%%%%%%%%%%%%%%%
\subsection{Can the Real Numbers Be Listed?}
\label{subsec:can-reals-be-listed}
%%%%%%%%%%%%%%%%%%%%%%%%%%%%%%%%%%%%%%%%%%%%%%%%%%%%%%%%%%%%

Suppose one attempts to list all real numbers in the interval

\[ (0,1). \]

If this were possible, the list might look like

\[
\begin{aligned}
x_1&=0.d_{11}d_{12}d_{13}d_{14}\ldots,\\
x_2&=0.d_{21}d_{22}d_{23}d_{24}\ldots,\\
x_3&=0.d_{31}d_{32}d_{33}d_{34}\ldots,\\
&\phantom{=}~\vdots
\end{aligned}
\]

where each \(d_{ij}\) is a decimal digit.

Cantor's diagonal argument shows that no such complete list can exist.

%%%%%%%%%%%%%%%%%%%%%%%%%%%%%%%%%%%%%%%%%%%%%%%%%%%%%%%%%%%%
\subsection{Cantor's Diagonal Argument}
\label{subsec:cantor-diagonal}
%%%%%%%%%%%%%%%%%%%%%%%%%%%%%%%%%%%%%%%%%%%%%%%%%%%%%%%%%%%%

Assume for contradiction that every real number in \((0,1)\) appears in a
list:

\[ x_1,x_2,x_3,\ldots. \]

Construct a new real number

\[ y=0.c_1c_2c_3\ldots \]

by choosing \(c_n\) so that it differs from the \(n\)-th decimal digit of
\(x_n\).

For example, one may define

\[ c_n=
\begin{cases}
1,&d_{nn}\neq1,\\
2,&d_{nn}=1.
\end{cases} \]

Then \(y\) differs from \(x_1\) in the first digit, from \(x_2\) in the second
digit, from \(x_3\) in the third digit, and so on.

Therefore

\[ y\neq x_n \]

for every \(n\).

But \(y\in(0,1)\), so the supposed list omitted a real number.

This is a contradiction.

%%%%%%%%%%%%%%%%%%%%%%%%%%%%%%%%%%%%%%%%%%%%%%%%%%%%%%%%%%%%
\subsection{The Real Numbers Are Uncountable}
\label{subsec:reals-uncountable}
%%%%%%%%%%%%%%%%%%%%%%%%%%%%%%%%%%%%%%%%%%%%%%%%%%%%%%%%%%%%

Cantor's argument proves that

\[ (0,1) \]

is uncountable.

Since

\[ (0,1)\subseteq\R, \]

the real numbers must also be uncountable.

Therefore,

\[ |\R|>\aleph_0. \]

There are strictly more real numbers than natural numbers.

%%%%%%%%%%%%%%%%%%%%%%%%%%%%%%%%%%%%%%%%%%%%%%%%%%%%%%%%%%%%
\subsection{Why the Diagonal Construction Works}
\label{subsec:why-diagonal-works}
%%%%%%%%%%%%%%%%%%%%%%%%%%%%%%%%%%%%%%%%%%%%%%%%%%%%%%%%%%%%

The construction does not merely produce a number different from many entries.

It produces a number different from \emph{every} entry.

For each \(n\),

\[ y \]

differs from

\[ x_n \]

specifically at the \(n\)-th decimal place.

Therefore no \(x_n\) can equal \(y\).

%%%%%%%%%%%%%%%%%%%%%%%%%%%%%%%%%%%%%%%%%%%%%%%%%%%%%%%%%%%%
\subsection{Decimal Representation Subtlety}
\label{subsec:decimal-representation-subtlety}
%%%%%%%%%%%%%%%%%%%%%%%%%%%%%%%%%%%%%%%%%%%%%%%%%%%%%%%%%%%%

Some real numbers have two decimal representations.

For example,

\[ 0.5000\ldots=0.4999\ldots. \]

When presenting a rigorous diagonal argument, one avoids this ambiguity by
choosing constructed digits from a restricted set, such as \(1\) and \(2\),
rather than using \(9\).

This ensures that the constructed decimal does not end in an infinite string
of \(9\)'s.

%%%%%%%%%%%%%%%%%%%%%%%%%%%%%%%%%%%%%%%%%%%%%%%%%%%%%%%%%%%%
\subsection{Not All Infinities Are Equal}
\label{subsec:not-all-infinities-equal}
%%%%%%%%%%%%%%%%%%%%%%%%%%%%%%%%%%%%%%%%%%%%%%%%%%%%%%%%%%%%

We have now seen

\[ |\N|=|\Z|=|\Q|=\aleph_0 \]

but

\[ |\R|>\aleph_0. \]

Thus there are different sizes of infinity.

This was one of Georg Cantor's most important discoveries.

%%%%%%%%%%%%%%%%%%%%%%%%%%%%%%%%%%%%%%%%%%%%%%%%%%%%%%%%%%%%
\subsection{The Continuum}
\label{subsec:continuum-cardinality}
%%%%%%%%%%%%%%%%%%%%%%%%%%%%%%%%%%%%%%%%%%%%%%%%%%%%%%%%%%%%

The cardinality of the real numbers is commonly denoted

\[ \mathfrak c. \]

The symbol

\[ \mathfrak c \]

is a lowercase \(c\) written in \emph{Fraktur} type and is read

\begin{center}
``the cardinality of the continuum''
\end{center}

or simply

\begin{center}
``the continuum.''
\end{center}

Thus,

\[ |\R|=\mathfrak c. \]

%%%%%%%%%%%%%%%%%%%%%%%%%%%%%%%%%%%%%%%%%%%%%%%%%%%%%%%%%%%%
\subsection{Intervals Have the Same Cardinality as \(\R\)}
\label{subsec:intervals-cardinality}
%%%%%%%%%%%%%%%%%%%%%%%%%%%%%%%%%%%%%%%%%%%%%%%%%%%%%%%%%%%%

A bounded interval can have the same cardinality as the entire real line.

For example,

\[ (0,1) \]

and

\[ \R \]

have the same cardinality.

One explicit bijection is

\[ f:(0,1)\to\R,\qquad
f(x)=\tan\left(\pi x-\frac{\pi}{2}\right). \]

As \(x\) moves through \((0,1)\), the argument of tangent moves through

\[ \left(-\frac{\pi}{2},\frac{\pi}{2}\right), \]

and the tangent function takes every real value exactly once.

Hence,

\[ |(0,1)|=|\R|=\mathfrak c. \]

%%%%%%%%%%%%%%%%%%%%%%%%%%%%%%%%%%%%%%%%%%%%%%%%%%%%%%%%%%%%
\subsection{Bounded Length Does Not Determine Cardinality}
\label{subsec:length-not-cardinality}
%%%%%%%%%%%%%%%%%%%%%%%%%%%%%%%%%%%%%%%%%%%%%%%%%%%%%%%%%%%%

The interval

\[ (0,1) \]

has finite geometric length \(1\).

The real line

\[ \R \]

has infinite geometric extent.

Yet they have the same cardinality.

Thus geometric size and set-theoretic cardinality measure fundamentally
different properties.

%%%%%%%%%%%%%%%%%%%%%%%%%%%%%%%%%%%%%%%%%%%%%%%%%%%%%%%%%%%%
\subsection{Open and Closed Intervals}
\label{subsec:open-closed-same-cardinality}
%%%%%%%%%%%%%%%%%%%%%%%%%%%%%%%%%%%%%%%%%%%%%%%%%%%%%%%%%%%%

For real numbers \(a<b\), each of the intervals

\[ (a,b),\qquad [a,b],\qquad [a,b),\qquad (a,b] \]

has cardinality

\[ \mathfrak c. \]

Adding or removing finitely many endpoints does not change the cardinality of
an infinite continuum-sized set.

%%%%%%%%%%%%%%%%%%%%%%%%%%%%%%%%%%%%%%%%%%%%%%%%%%%%%%%%%%%%
\subsection{Comparing Cardinalities}
\label{subsec:comparing-cardinalities}
%%%%%%%%%%%%%%%%%%%%%%%%%%%%%%%%%%%%%%%%%%%%%%%%%%%%%%%%%%%%

The notation

\[ |A|\leq|B| \]

can be interpreted as saying that there exists an injective function

\[ f:A\to B. \]

This means that the elements of \(A\) can be embedded into \(B\) without
collisions.

If

\[ |A|\leq|B| \]

and

\[ |B|\leq|A|, \]

one might expect that

\[ |A|=|B|. \]

This is indeed true.

%%%%%%%%%%%%%%%%%%%%%%%%%%%%%%%%%%%%%%%%%%%%%%%%%%%%%%%%%%%%
\subsection{Cantor--Bernstein--Schröder Theorem}
\label{subsec:cantor-bernstein-schroeder}
%%%%%%%%%%%%%%%%%%%%%%%%%%%%%%%%%%%%%%%%%%%%%%%%%%%%%%%%%%%%

\begin{theorem}[Cantor--Bernstein--Schröder]
If there exists an injection
\[ f:A\to B \]
and an injection
\[ g:B\to A, \]
then there exists a bijection
\[ h:A\to B. \]

Therefore,
\[ |A|=|B|. \]
\end{theorem}

The theorem is also commonly called the Schröder--Bernstein theorem.

It is especially useful when constructing a direct bijection would be
difficult.

%%%%%%%%%%%%%%%%%%%%%%%%%%%%%%%%%%%%%%%%%%%%%%%%%%%%%%%%%%%%
\subsection{Using Cantor--Bernstein--Schröder}
\label{subsec:using-cbs}
%%%%%%%%%%%%%%%%%%%%%%%%%%%%%%%%%%%%%%%%%%%%%%%%%%%%%%%%%%%%

Suppose one wants to show

\[ |A|=|B|. \]

Instead of explicitly constructing a bijection, it may be easier to construct:

\[ A\hookrightarrow B \]

and

\[ B\hookrightarrow A. \]

The hooked arrow

\[ \hookrightarrow \]

is commonly read

\begin{center}
``injects into''
\end{center}

and is used informally to indicate an injective mapping.

The Cantor--Bernstein--Schröder theorem then guarantees equal cardinality.

%%%%%%%%%%%%%%%%%%%%%%%%%%%%%%%%%%%%%%%%%%%%%%%%%%%%%%%%%%%%
\subsection{Power Sets Revisited}
\label{subsec:power-sets-cardinality}
%%%%%%%%%%%%%%%%%%%%%%%%%%%%%%%%%%%%%%%%%%%%%%%%%%%%%%%%%%%%

Recall that

\[ \mathcal P(A) \]

is the set of all subsets of \(A\).

For finite sets,

\[ |\mathcal P(A)|=2^{|A|}. \]

For infinite sets, power sets also have strictly greater cardinality than the
original set.

This is Cantor's theorem.

%%%%%%%%%%%%%%%%%%%%%%%%%%%%%%%%%%%%%%%%%%%%%%%%%%%%%%%%%%%%
\subsection{Cantor's Theorem}
\label{subsec:cantor-theorem}
%%%%%%%%%%%%%%%%%%%%%%%%%%%%%%%%%%%%%%%%%%%%%%%%%%%%%%%%%%%%

\begin{theorem}[Cantor's Theorem]
For every set \(A\),
\[ |A|<|\mathcal P(A)|. \]
\end{theorem}

In other words, there is no bijection from a set to its own power set.

This theorem is true for both finite and infinite sets.

%%%%%%%%%%%%%%%%%%%%%%%%%%%%%%%%%%%%%%%%%%%%%%%%%%%%%%%%%%%%
\subsection{The Diagonal Set in Cantor's Theorem}
\label{subsec:cantor-diagonal-set}
%%%%%%%%%%%%%%%%%%%%%%%%%%%%%%%%%%%%%%%%%%%%%%%%%%%%%%%%%%%%

Suppose for contradiction that there were a surjective function

\[ f:A\to\mathcal P(A). \]

Construct the set

\[ D=\{x\in A:x\notin f(x)\}. \]

Since \(D\subseteq A\),

\[ D\in\mathcal P(A). \]

If \(f\) were surjective, there would exist some \(d\in A\) such that

\[ f(d)=D. \]

Now ask whether

\[ d\in D. \]

By definition of \(D\),

\[ d\in D\Longleftrightarrow d\notin f(d). \]

But \(f(d)=D\), so

\[ d\in D\Longleftrightarrow d\notin D. \]

This is impossible.

Therefore no such surjection exists.

%%%%%%%%%%%%%%%%%%%%%%%%%%%%%%%%%%%%%%%%%%%%%%%%%%%%%%%%%%%%
\subsection{Worked Example: Why No Set Matches Its Power Set}
\label{subsec:worked-cantor-theorem}
%%%%%%%%%%%%%%%%%%%%%%%%%%%%%%%%%%%%%%%%%%%%%%%%%%%%%%%%%%%%

\begin{workedexamplebox}{Finite Preview of Cantor's Theorem}

Let

\[ A=\{a,b,c\}. \]

Then

\[ |A|=3. \]

Its power set contains

\[ 2^3=8 \]

elements.

Thus,

\begin{tabularx}{\textwidth}{@{}>{\raggedright\arraybackslash}p{0.42\textwidth}X@{}}
Original set size. & \(|A|=3\)\\
Power set size. & \(|\mathcal P(A)|=8\)\\
Compare. & \(3<8\)
\end{tabularx}

\begin{answerbox}
\[ |A|<|\mathcal P(A)| \]
\end{answerbox}
\end{workedexamplebox}

%%%%%%%%%%%%%%%%%%%%%%%%%%%%%%%%%%%%%%%%%%%%%%%%%%%%%%%%%%%%
\subsection{The Power Set of \(\N\)}
\label{subsec:power-set-n}
%%%%%%%%%%%%%%%%%%%%%%%%%%%%%%%%%%%%%%%%%%%%%%%%%%%%%%%%%%%%

Cantor's theorem gives

\[ |\N|<|\mathcal P(\N)|. \]

In fact,

\[ |\mathcal P(\N)|=\mathfrak c. \]

Thus the collection of all subsets of the natural numbers has the same
cardinality as the real numbers.

This is a remarkable connection between discrete and continuous mathematics.

%%%%%%%%%%%%%%%%%%%%%%%%%%%%%%%%%%%%%%%%%%%%%%%%%%%%%%%%%%%%
\subsection{Binary Sequences and Subsets of \(\N\)}
\label{subsec:binary-sequences-subsets}
%%%%%%%%%%%%%%%%%%%%%%%%%%%%%%%%%%%%%%%%%%%%%%%%%%%%%%%%%%%%

Every subset

\[ S\subseteq\N \]

can be represented by an infinite binary sequence.

Write

\[ a_n=
\begin{cases}
1,&n\in S,\\
0,&n\notin S.
\end{cases} \]

The sequence

\[ a_1a_2a_3\ldots \]

records exactly which natural numbers belong to \(S\).

Thus subsets of \(\N\) correspond naturally to infinite sequences of zeros
and ones.

%%%%%%%%%%%%%%%%%%%%%%%%%%%%%%%%%%%%%%%%%%%%%%%%%%%%%%%%%%%%
\subsection{No Largest Infinity}
\label{subsec:no-largest-infinity}
%%%%%%%%%%%%%%%%%%%%%%%%%%%%%%%%%%%%%%%%%%%%%%%%%%%%%%%%%%%%

Cantor's theorem has an extraordinary consequence.

Given any set \(A\),

\[ |\mathcal P(A)|>|A|. \]

Then

\[ |\mathcal P(\mathcal P(A))|>|\mathcal P(A)|, \]

and continuing gives

\[ |A|
<
|\mathcal P(A)|
<
|\mathcal P(\mathcal P(A))|
<
\cdots. \]

Therefore there is no largest cardinality.

For every size of infinity, there is a strictly larger one.

%%%%%%%%%%%%%%%%%%%%%%%%%%%%%%%%%%%%%%%%%%%%%%%%%%%%%%%%%%%%
\subsection{Infinity Is Not an Ordinary Number}
\label{subsec:infinity-not-ordinary-number}
%%%%%%%%%%%%%%%%%%%%%%%%%%%%%%%%%%%%%%%%%%%%%%%%%%%%%%%%%%%%

The symbol

\[ \infty \]

is called the \emph{infinity symbol} and is read ``infinity.''

In introductory calculus and interval notation, \(\infty\) often describes
unbounded behavior.

It should not automatically be treated as an ordinary real number.

Cardinal numbers such as

\[ \aleph_0 \]

are mathematical objects used specifically to measure infinite cardinality.

Thus the symbols

\[ \infty,\qquad \aleph_0,\qquad \mathfrak c \]

have related but distinct roles.

%%%%%%%%%%%%%%%%%%%%%%%%%%%%%%%%%%%%%%%%%%%%%%%%%%%%%%%%%%%%
\subsection{Cardinality Does Not Measure Density}
\label{subsec:cardinality-not-density}
%%%%%%%%%%%%%%%%%%%%%%%%%%%%%%%%%%%%%%%%%%%%%%%%%%%%%%%%%%%%

Consider

\[ \N\subset\Q\subset\R. \]

The natural numbers are sparse on the real line.

The rationals are dense.

The reals fill the continuum.

Yet

\[ |\N|=|\Q| \]

while

\[ |\R|>|\Q|. \]

Spatial distribution and cardinality are different mathematical concepts.

%%%%%%%%%%%%%%%%%%%%%%%%%%%%%%%%%%%%%%%%%%%%%%%%%%%%%%%%%%%%
\subsection{Cardinality Does Not Measure Geometric Length}
\label{subsec:cardinality-not-length}
%%%%%%%%%%%%%%%%%%%%%%%%%%%%%%%%%%%%%%%%%%%%%%%%%%%%%%%%%%%%

The intervals

\[ (0,1) \]

and

\[ (0,1000000) \]

have dramatically different lengths.

Nevertheless,

\[ |(0,1)|=|(0,1000000)|=\mathfrak c. \]

A simple bijection is

\[ f(x)=1000000x. \]

Thus equal cardinality does not imply equal geometric measure.

%%%%%%%%%%%%%%%%%%%%%%%%%%%%%%%%%%%%%%%%%%%%%%%%%%%%%%%%%%%%
\subsection{Removing Finitely Many Elements}
\label{subsec:removing-finite-elements}
%%%%%%%%%%%%%%%%%%%%%%%%%%%%%%%%%%%%%%%%%%%%%%%%%%%%%%%%%%%%

Removing finitely many elements from a countably infinite set does not change
its cardinality.

For example,

\[ \N\setminus\{1,2,3,4,5\} \]

is still countably infinite.

Similarly,

\[ \R\setminus\{0\} \]

still has cardinality

\[ \mathfrak c. \]

%%%%%%%%%%%%%%%%%%%%%%%%%%%%%%%%%%%%%%%%%%%%%%%%%%%%%%%%%%%%
\subsection{Adding Countably Many Elements}
\label{subsec:adding-countably-many}
%%%%%%%%%%%%%%%%%%%%%%%%%%%%%%%%%%%%%%%%%%%%%%%%%%%%%%%%%%%%

If \(A\) has cardinality \(\mathfrak c\) and \(B\) is countable, then in
standard situations

\[ |A\cup B|=\mathfrak c. \]

For example,

\[ |\R\cup\N|=|\R| \]

because

\[ \N\subseteq\R. \]

The larger infinite cardinality absorbs the smaller one in this union.

%%%%%%%%%%%%%%%%%%%%%%%%%%%%%%%%%%%%%%%%%%%%%%%%%%%%%%%%%%%%
\subsection{Common Cardinality Errors}
\label{subsec:common-cardinality-errors}
%%%%%%%%%%%%%%%%%%%%%%%%%%%%%%%%%%%%%%%%%%%%%%%%%%%%%%%%%%%%

\subsubsection{Assuming Proper Subsets Are Always Smaller}

For finite sets,

\[ A\subsetneq B\Longrightarrow|A|<|B|. \]

For infinite sets, this implication can fail.

For example,

\[ 2\N\subsetneq\N \]

but

\[ |2\N|=|\N|. \]

%%%%%%%%%%%%%%%%%%%%%%%%%%%%%%%%%%%%%%%%%%%%%%%%%%%%%%%%%%%%
\subsubsection{Assuming \(\Z\) Is Twice as Large as \(\N\)}
\label{subsec:error-z-twice-n}
%%%%%%%%%%%%%%%%%%%%%%%%%%%%%%%%%%%%%%%%%%%%%%%%%%%%%%%%%%%%

The integers contain positive and negative values, but cardinality does not
behave like ordinary finite counting.

There exists a bijection

\[ \N\to\Z. \]

Therefore,

\[ |\N|=|\Z|. \]

%%%%%%%%%%%%%%%%%%%%%%%%%%%%%%%%%%%%%%%%%%%%%%%%%%%%%%%%%%%%
\subsubsection{Assuming Dense Means Uncountable}
\label{subsec:error-dense-uncountable}
%%%%%%%%%%%%%%%%%%%%%%%%%%%%%%%%%%%%%%%%%%%%%%%%%%%%%%%%%%%%

The rational numbers are dense in \(\R\), but

\[ |\Q|=\aleph_0. \]

Therefore a dense subset of the real line may still be countable.

%%%%%%%%%%%%%%%%%%%%%%%%%%%%%%%%%%%%%%%%%%%%%%%%%%%%%%%%%%%%
\subsubsection{Assuming Every Infinite Set Is Countable}
\label{subsec:error-all-infinite-countable}
%%%%%%%%%%%%%%%%%%%%%%%%%%%%%%%%%%%%%%%%%%%%%%%%%%%%%%%%%%%%

The real numbers are infinite but uncountable.

Thus

\[ \aleph_0<\mathfrak c. \]

There are multiple sizes of infinity.

%%%%%%%%%%%%%%%%%%%%%%%%%%%%%%%%%%%%%%%%%%%%%%%%%%%%%%%%%%%%
\subsubsection{Treating Infinity Like a Real Number}
\label{subsec:error-infinity-real}
%%%%%%%%%%%%%%%%%%%%%%%%%%%%%%%%%%%%%%%%%%%%%%%%%%%%%%%%%%%%

Expressions such as

\[ \infty+1=\infty \]

may appear informally in certain mathematical settings, but \(\infty\) is not
an ordinary real number.

Infinite cardinal arithmetic and extended-real arithmetic require their own
definitions and should not be inferred from ordinary arithmetic rules.

%%%%%%%%%%%%%%%%%%%%%%%%%%%%%%%%%%%%%%%%%%%%%%%%%%%%%%%%%%%%
\subsection{Worked Examples}
\label{subsec:cardinality-worked-examples}
%%%%%%%%%%%%%%%%%%%%%%%%%%%%%%%%%%%%%%%%%%%%%%%%%%%%%%%%%%%%

\begin{workedexamplebox}{Show the Odd Positive Integers Are Countable}

Let

\[ O=\{1,3,5,7,\ldots\}. \]

Define

\[ f:\N\to O,\qquad f(n)=2n-1. \]

\begin{tabularx}{\textwidth}{@{}>{\raggedright\arraybackslash}p{0.42\textwidth}X@{}}
Every \(n\in\N\) produces an odd positive integer. & \(f(n)=2n-1\)\\
Equal outputs imply equal inputs. & \(2m-1=2n-1\Rightarrow m=n\)\\
Every odd positive integer has form \(2n-1\). & \(f\) is surjective
\end{tabularx}

\begin{answerbox}
\[ |O|=|\N|=\aleph_0 \]
\end{answerbox}
\end{workedexamplebox}

\begin{workedexamplebox}{Compare \((0,1)\) and \((0,5)\)}

Define

\[ f:(0,1)\to(0,5),\qquad f(x)=5x. \]

\begin{tabularx}{\textwidth}{@{}>{\raggedright\arraybackslash}p{0.42\textwidth}X@{}}
If \(x\in(0,1)\). & \(5x\in(0,5)\)\\
Equal outputs imply equal inputs. & \(5x_1=5x_2\Rightarrow x_1=x_2\)\\
Given \(y\in(0,5)\). & \(x=y/5\in(0,1)\)
\end{tabularx}

Thus \(f\) is bijective.

\begin{answerbox}
\[ |(0,1)|=|(0,5)|=\mathfrak c \]
\end{answerbox}
\end{workedexamplebox}

\begin{workedexamplebox}{Countability of Integer Pairs}

Consider

\[ \Z\times\Z. \]

Since

\[ |\Z|=|\N|, \]

and

\[ |\N\times\N|=|\N|, \]

the product of two countable sets is countable.

\begin{answerbox}
\[ |\Z\times\Z|=\aleph_0 \]
\end{answerbox}
\end{workedexamplebox}

%%%%%%%%%%%%%%%%%%%%%%%%%%%%%%%%%%%%%%%%%%%%%%%%%%%%%%%%%%%%
\subsection{Exercises}
\label{subsec:cardinality-exercises}
%%%%%%%%%%%%%%%%%%%%%%%%%%%%%%%%%%%%%%%%%%%%%%%%%%%%%%%%%%%%

\begin{exercise}[1]
Find the cardinality of
\[ A=\{2,4,6,8,10,12\}. \]
\end{exercise}

\begin{exercise}[1]
Find
\[ |\mathcal P(\{a,b,c,d\})|. \]
\end{exercise}

\begin{exercise}[1]
Construct a bijection between
\[ \N \]
and
\[ \{5,10,15,20,\ldots\}. \]
\end{exercise}

\begin{exercise}[1]
Construct a bijection between \(\N\) and the positive odd integers.
\end{exercise}

\begin{exercise}[2]
Explain why
\[ \N\setminus\{1\} \]
has the same cardinality as \(\N\).
\end{exercise}

\begin{exercise}[2]
Construct an explicit bijection between
\[ \N \]
and
\[ \N\setminus\{1,2,3\}. \]
\end{exercise}

\begin{exercise}[2]
Explain why the integers are countable.
\end{exercise}

\begin{exercise}[2]
Describe a systematic listing of
\[ \N\times\N. \]
\end{exercise}

\begin{exercise}[2]
Explain why every infinite subset of \(\N\) is countably infinite.
\end{exercise}

\begin{exercise}[2]
Explain why the set of even integers is countable.
\end{exercise}

\begin{exercise}[2]
Explain why the set of rational numbers is countable.
\end{exercise}

\begin{exercise}[2]
Explain why density does not imply uncountability.
\end{exercise}

\begin{exercise}[2]
Give an example of a countable dense subset of \(\R\).
\end{exercise}

\begin{exercise}[2]
Give an example of an uncountable subset of \(\R\).
\end{exercise}

\begin{exercise}[2]
Explain why
\[ |\N|<|\mathcal P(\N)|. \]
\end{exercise}

\begin{exercise}[2]
If
\[ |A|=6, \]
find
\[ |\mathcal P(A)|. \]
\end{exercise}

\begin{exercise}[2]
Construct a bijection between
\[ (0,1) \]
and
\[ (2,8). \]
\end{exercise}

\begin{exercise}[2]
Construct a bijection between
\[ \R \]
and
\[ (0,\infty). \]
\end{exercise}

\begin{exercise}[3]
Show that
\[ |\N\times\{1,2,3\}|=\aleph_0. \]
\end{exercise}

\begin{exercise}[3]
Show that
\[ \Z\times\N \]
is countable.
\end{exercise}

\begin{exercise}[3]
Prove that a countable union of finite sets is countable.
\end{exercise}

\begin{exercise}[3]
Prove that a finite union of countable sets is countable.
\end{exercise}

\begin{exercise}[3]
Explain carefully why Cantor's diagonal number cannot appear anywhere in the
supposed enumeration of \((0,1)\).
\end{exercise}

\begin{exercise}[3]
Prove that
\[ \R\setminus\{0\} \]
has the same cardinality as \(\R\).
\end{exercise}

\begin{exercise}[3]
Use the Cantor--Bernstein--Schröder theorem to show that
\[ |[0,1]|=|(0,1)|. \]
\end{exercise}

\begin{exercise}[3]
Suppose \(A\) is any set. Explain the role of
\[ D=\{x\in A:x\notin f(x)\} \]
in the proof of Cantor's theorem.
\end{exercise}

\begin{exercise}[3]
Explain why Cantor's theorem implies that no largest cardinality exists.
\end{exercise}

%%%%%%%%%%%%%%%%%%%%%%%%%%%%%%%%%%%%%%%%%%%%%%%%%%%%%%%%%%%%
\subsection{Section Connections}
\label{subsec:cardinality-connections}
%%%%%%%%%%%%%%%%%%%%%%%%%%%%%%%%%%%%%%%%%%%%%%%%%%%%%%%%%%%%

\chapterconnection{
Cardinality combines set theory, functions, logic, and proof methods from the
earlier sections of Chapter 0. Countability becomes important in number theory,
discrete mathematics, probability, topology, and analysis. Uncountability is
central to the structure of the real numbers and later motivates deeper
questions in axiomatic set theory. Power-set arguments and Cantor's theorem
also reappear in mathematical logic and foundations.
}

%%%%%%%%%%%%%%%%%%%%%%%%%%%%%%%%%%%%%%%%%%%%%%%%%%%%%%%%%%%%
% SECTION SOURCE TRACKING
%%%%%%%%%%%%%%%%%%%%%%%%%%%%%%%%%%%%%%%%%%%%%%%%%%%%%%%%%%%%

%------------------------------------------------------------
% 0.9 Cardinality and Infinity
%------------------------------------------------------------
%
% Sources:
%   - Richard Hammack, Book of Proof, Third Edition.
%     Key: hammack2018bookofproof
%
% Used for:
%   - cardinality through bijections;
%   - countable and countably infinite sets;
%   - countability of Z and Q;
%   - diagonal enumeration arguments;
%   - Cantor's diagonal argument;
%   - uncountability of the real numbers;
%   - comparison of cardinalities;
%   - Cantor--Bernstein--Schroeder theorem;
%   - Cantor's theorem and power sets.
%
% Notes:
%   - Examples and exposition are independently developed.
%   - Formal cardinal arithmetic and axiomatic set theory are deferred to
%     Chapter 12.
%   - Measure, density, and continuum structure are developed more deeply in
%     Analysis.
%
\nocite{hammack2018bookofproof}